\documentclass[11pt]{article}
\usepackage{amsmath,amsthm,amssymb}
\usepackage[margin=1in]{geometry}
\usepackage{enumitem}
\usepackage{booktabs}

\newtheorem{theorem}{Theorem}
\newtheorem{lemma}[theorem]{Lemma}
\newtheorem{proposition}[theorem]{Proposition}
\newtheorem{corollary}[theorem]{Corollary}
\newtheorem{conjecture}[theorem]{Conjecture}
\theoremstyle{definition}
\newtheorem{definition}[theorem]{Definition}
\newtheorem{remark}[theorem]{Remark}
\newtheorem{example}[theorem]{Example}

\title{Uniform per-trajectory invariant for the\\
  $(2^a,1^{n-2a})$ $j$-formula support flow}
\author{Clio}
\date{17 May 2026}

\begin{document}
\maketitle

\begin{abstract}
We give a uniform structural proof of the \emph{per-trajectory
invariant} (Conjecture~4.4 of \texttt{2026-05-16-a4-support-flow.tex})
for the family $\lambda=(2^a,1^{n-2a})$ at every $a\ge 2$ and
$n\ge 3a$.  This invariant was previously proved by case-by-case
enumeration at $a=2,3,4$ (May-13, May-15-a3, May-16); the May-16
paper left the general-$a$ case as a conjecture requiring either
$2^{a-1}$ cases per step at each fixed $a$ or a uniform positional
invariant.  We give the uniform proof.

The argument reduces to a four-case structural lemma that depends
only on \emph{which columns of $\lambda$} the three consecutive
entries $\{n-2-k,n-1-k,n-k\}$ inhabit, not on the encoding bits
$\tau\in\{L,R\}^{a-1}$.  The case bound $n\ge 3a$ enters once, via
the same-row exclusion (May-16 Claim~(4)).

\textbf{Consequence.}  Combined with the May-15 inductive reduction
(\texttt{2026-05-15-general-a-inductive.tex}, Theorem~3.2) and the
May-14 Phase~A endpoint at general $\lambda$ (\texttt{2026-05-14-2a-1n-family.tex},
Proposition~3.1), the rank-zero theorem
$\Pi^{S_n}|_{V_{(2^a,1^{n-2a})}}=0$ holds unconditionally for every
$a\ge 2$ and $n\ge 3a$.  In particular, the per-trajectory
framework now yields the rank-zero result for all of the family
$\lambda=(2^a,1^{n-2a})$ simultaneously, modulo only the (mechanical)
general-$a$ Phase~A.

The support-size $2^a$ claim and the L/R encoding identification of
Conjecture~4.5 (May-16) remain open as written; we comment on the
finite combinatorial residue at the end.
\end{abstract}

\section{Setup}\label{sec:setup}

We adopt the conventions of \texttt{2026-05-15-a3-via-branching.tex},
\texttt{2026-05-15-general-a-inductive.tex}, and
\texttt{2026-05-16-a4-support-flow.tex} (which we abbreviate
\textbf{May-15-a3}, \textbf{May-15-gen}, and \textbf{May-16}).
$H_q(S_n)$ is the Hecke algebra over $\mathbb{Q}(q)$ with generators
$T_1,\ldots,T_{n-1}$.  For $\lambda\vdash n$, $V_\lambda$ is the
seminormal cell module with basis $\{v_T\}$ indexed by SYTs of
$\lambda$.  $E_j^\pm$ is the $\pm$-eigenspace of $T_j$.

Throughout, fix $a\ge 2$ and $n\ge 3a$:
\[
   \lambda:=(2^a,1^{n-2a})\vdash n,\qquad
   \mu_1:=(2^{a-1},1^{n-2a})\vdash n-2.
\]
$\ell(\lambda)=n-a$.  Define
\[
   R'_k=(T_{k-1}{+}1)\cdots(T_1{+}1),\qquad
   B_j=R'_jR'_{j+1}\cdots R'_n.
\]

Per May-15-gen Theorem~3.2 (the inductive reduction),
\begin{equation}\label{eq:gen-reduction}
   B_{n-a+1}^{(\lambda)}V_\lambda
   = (T_{n-a}{+}1)(T_{n-a+1}{+}1)\cdots(T_{n-2}{+}1)\,
   \Phi_1\big(B_{n-a}^{(\mu_1)}V_{\mu_1}\big),
\end{equation}
where $\Phi_1:V_{\mu_1}|_{n-2}\xrightarrow{\sim}E_{n-1}^+|_{V_\lambda|_n}$
is the branching isomorphism of May-15-gen Definition~3.1:
for each SYT $T'$ of $\mu_1$, $\Phi_1(v_{T'})$ is the
$T_{n-1}$-$+q$-eigenvector in the 2-block
$\{v_{T_A(T')},v_{T_B(T')}\}$, where
\begin{itemize}[leftmargin=2em,topsep=0.3em,itemsep=0.2em]
\item $T_A(T')$: extend $T'$ by placing $n-1$ at $(a,2)$, $n$ at
$(n-a,1)$;
\item $T_B(T')$: extend $T'$ by placing $n$ at $(a,2)$, $n-1$ at
$(n-a,1)$.
\end{itemize}

\subsection{Trajectory}

Following May-16, for each SYT $T'$ of $\mu_1$ define
\[
   w_k(T')\;:=\;(T_{n-1-k}{+}1)(T_{n-k}{+}1)\cdots(T_{n-2}{+}1)\,
   \Phi_1(v_{T'}^{(\mu_1)}),\qquad k=0,1,\ldots,a-1.
\]
So $w_0(T')=\Phi_1(v_{T'})$ and $w_k(T')=(T_{n-1-k}{+}1)\,w_{k-1}(T')$.

\begin{definition}[Per-trajectory invariant]\label{def:invariant}
We say $T'$ satisfies the \emph{per-trajectory invariant at step $k$}
($k\in\{0,1,\ldots,a-1\}$) if:
\begin{enumerate}[leftmargin=2em,topsep=0.3em,itemsep=0.2em,label=(P\arabic*)]
\item $w_k(T')\ne 0$.
\item Its seminormal support is exactly two SYTs, say
$\{P_k^+,P_k^-\}\subseteq\mathrm{SYT}(\lambda)$.
\item $P_k^-=s_{n-1-k}P_k^+$ (they form a $T_{n-1-k}$-2-block),
and the 2-block is non-degenerate ($|c(n-k)-c(n-1-k)|\ge 2$ in $P_k^+$).
\item Both coefficients of $w_k(T')$ in the seminormal basis are
nonzero in $\mathbb{Q}(q)$.
\end{enumerate}
\end{definition}

The base case $k=0$ is the definition of $\Phi_1$.

\section{Main result}\label{sec:main}

\begin{theorem}[Uniform per-trajectory invariant]\label{thm:main}
For every $a\ge 2$, every $n\ge 3a$, every SYT $T'$ of $\mu_1$, and
every $k\in\{0,1,\ldots,a-1\}$, the per-trajectory invariant
(Definition~\ref{def:invariant}) holds.
\end{theorem}

This is May-16 Conjecture~4.4, proved here without case enumeration.
The proof is induction on $k$ from the structural lemma below.

\begin{lemma}[Kill/expand alternation]\label{lem:kill-expand}
Let $k\in\{0,1,\ldots,a-2\}$ and suppose
$\{P_k^+,P_k^-=s_{n-1-k}P_k^+\}\subseteq\mathrm{SYT}(\lambda)$ is a
non-degenerate $T_{n-1-k}$-2-block.  Set
\[
   e_1:=n-2-k,\quad e_2:=n-1-k,\quad e_3:=n-k.
\]
Then exactly one of $\{P_k^+,P_k^-\}$ is annihilated by
$(T_{e_1}{+}1)$, and the other yields a non-degenerate
$T_{e_1}$-2-block via the Hoefsmit action.
\end{lemma}

The base case $k=0$ of Theorem~\ref{thm:main} is immediate from the
construction of $\Phi_1$ (May-15-gen Definition~3.1, Lemma~3.4):
$w_0(T')$ is by construction the $T_{n-1}$-$+q$-eigenvector in the
2-block $\{v_{T_A(T')},v_{T_B(T')}\}$ with axial distance
$|c(n)-c(n-1)|=n-2a+1\ge a+1\ge 3$, non-degenerate.  All four
properties (P1)--(P4) hold.

The inductive step $k\to k+1$ of Theorem~\ref{thm:main} from
Lemma~\ref{lem:kill-expand}: given (P1)--(P4) at step $k$, write
$w_k(T')=\alpha v_{P_k^+}+\beta v_{P_k^-}$ with $\alpha,\beta\ne 0$.
By Lemma~\ref{lem:kill-expand}, exactly one of $\{v_{P_k^+},v_{P_k^-}\}$
is killed by $(T_{e_1}{+}1)$; say $v_{P_k^+}$ is killed.  Then
\[
   w_{k+1}(T')=(T_{e_1}{+}1)w_k(T')=\beta\cdot(T_{e_1}{+}1)v_{P_k^-}.
\]
By the same lemma, $(T_{e_1}{+}1)v_{P_k^-}$ is a Hoefsmit
$+q$-eigenvector in the non-degenerate $T_{e_1}$-2-block
$\{v_{P_k^-},v_{s_{e_1}P_k^-}\}$, hence nonzero with both coefficients
nonzero (the Hoefsmit eigenvector formula).  This gives
(P1)--(P4) at step $k+1$ with $P_{k+1}^+=P_k^-$ and
$P_{k+1}^-=s_{e_1}P_k^-$.  (The case $v_{P_k^-}$ killed is symmetric.)

\medskip
\noindent
So Theorem~\ref{thm:main} reduces to Lemma~\ref{lem:kill-expand},
which we now prove.

\section{Proof of Lemma~\ref{lem:kill-expand}}\label{sec:proof-lemma}

\subsection{Column structure of large entries}\label{sec:claim4}

We first establish that $e_1,e_2,e_3$ lie in one of only two
``allowed regions'' of $\lambda$.

\begin{lemma}[Column-row bound for large entries]\label{lem:bound}
In any SYT $T$ of $\lambda=(2^a,1^{n-2a})$, the column-$1$ row-$r$
entry for $r\le a$ is at most $2a-1$.
\end{lemma}

\begin{proof}
Among the entries $\{1,\ldots,2a-1\}$ ($2a-1$ values), at most $a-1$
can be in column $2$ (since the SYT condition forces column-$2$
row-$a$ to be strictly greater than column-$1$ row-$a$, so any
column-$2$ entry $\le 2a-1$ must lie in some row $r'\le a-1$).
Hence at least $a$ of these values lie in column $1$.  By
column-monotonicity, the $a$-th smallest column-$1$ entry is
$\le 2a-1$.  But that $a$-th smallest is precisely column-$1$ row $a$.
For $r<a$, the column-$1$ row-$r$ entry is even smaller.
\end{proof}

\begin{corollary}[Same-row exclusion]\label{cor:no-same-row}
For $k\in\{0,1,\ldots,a-2\}$ and any SYT $T$ of $\lambda$,
$\{e_1,e_2,e_3\}=\{n-2-k,n-1-k,n-k\}$ each lies in:
\begin{itemize}[leftmargin=2em,topsep=0.3em,itemsep=0.2em]
\item Column $1$ at some row $r>a$, or
\item Column $2$ at some row $r\le a$.
\end{itemize}
\end{corollary}

\begin{proof}
$e_1=n-2-k\ge n-2-(a-2)=n-a\ge 2a$ (using $n\ge 3a$).  By
Lemma~\ref{lem:bound}, $e_1$ cannot lie at column-$1$ row $\le a$.
Column-$2$ has no rows above $a$ (only $a$ rows in column $2$ of
$\lambda$).  So column-$1$ row $>a$ or column-$2$ row $\le a$.
The same applies to $e_2=e_1+1$, $e_3=e_1+2$.
\end{proof}

In particular, $e_1,e_2,e_3$ \emph{cannot share a length-$2$ row}
of $\lambda$ (which would require one in column-$1$ at row $\le a$).

\subsection{The four-case structural argument}\label{sec:four-cases}

We now prove Lemma~\ref{lem:kill-expand}.  The hypothesis is that
$\{P_k^+,P_k^-\}$ is a non-degenerate $T_{e_2}$-2-block.  This means:
in $P_k^+$, entries $e_2$ and $e_3$ lie at different rows and
different columns, with axial distance $\ge 2$.

By Corollary~\ref{cor:no-same-row}, each of $e_2,e_3$ is in
\{col-$1$ row $>a$\} $\cup$ \{col-$2$ row $\le a$\}; the two
``regions'' are row-disjoint.  Thus $e_2$ and $e_3$ occupy
\emph{different columns} (one in col-$1$, the other in col-$2$) and
automatically different rows.  Two sub-cases:

\paragraph{Case (I).}  In $P_k^+$: $e_2$ in column $1$ at row
$r_2>a$, $e_3$ in column $2$ at row $s_3\le a$.

In $P_k^-=s_{e_2}P_k^+$: $e_2$ in column $2$ at row $s_3$, $e_3$ in
column $1$ at row $r_2$.

\paragraph{Case (II).}  In $P_k^+$: $e_2$ in column $2$ at row
$s_2\le a$, $e_3$ in column $1$ at row $r_3>a$.

In $P_k^-$: $e_2$ in column $1$ at row $r_3$, $e_3$ in column $2$ at
row $s_2$.

\medskip

\noindent
By Corollary~\ref{cor:no-same-row}, $e_1$ also lies in one of the
two regions.  Two sub-cases each:

\subsubsection*{Case (I.a):  $e_1$ in column $1$ at row $r_1>a$.}

\begin{itemize}[leftmargin=2em,topsep=0.3em,itemsep=0.2em]
\item In $P_k^+$: $e_1,e_2$ both in column $1$.  No integer between
$e_1$ and $e_2$ ($e_2=e_1+1$), so no column-$1$ entry of $P_k^+$ lies
strictly between them; hence $e_1,e_2$ occupy \emph{consecutive
rows} of column $1$.  Same column, adjacent rows: $T_{e_1}$ has
eigenvalue $-1$ on $v_{P_k^+}$, so $(T_{e_1}{+}1)v_{P_k^+}=0$.
\textbf{Killed.}
\item In $P_k^-$: $e_1$ in column $1$, $e_2$ in column $2$.
Different columns.  $e_1$ at row $r_1>a$, $e_2$ at row $s_3\le a$:
different rows.  Hence $T_{e_1}$ acts on $\{v_{P_k^-},v_{s_{e_1}P_k^-}\}$
as a Hoefsmit 2-block.  Axial distance:
$|c(e_2)-c(e_1)|=|(2-s_3)-(1-r_1)|=|1+r_1-s_3|=1+r_1-s_3\ge 2$
(since $r_1\ge a+1$, $s_3\le a$).  Non-degenerate.  \textbf{Expanded.}
\end{itemize}

\subsubsection*{Case (I.b): $e_1$ in column $2$ at row $s_1\le a$.}

\begin{itemize}[leftmargin=2em,topsep=0.3em,itemsep=0.2em]
\item In $P_k^+$: $e_1$ in column $2$, $e_2$ in column $1$.
Different columns.  Different rows ($s_1\le a$, $r_2>a$).  Axial
distance $|c(e_2)-c(e_1)|=|(1-r_2)-(2-s_1)|=|s_1-1-r_2|=r_2+1-s_1\ge 2$.
Non-degenerate.  \textbf{Expanded.}
\item In $P_k^-$: $e_1$ in column $2$, $e_2$ in column $2$.  Same
column.  No integer between $e_1,e_2$, so they occupy consecutive
rows of column $2$ in $P_k^-$.  Adjacent column, same column:
$T_{e_1}v_{P_k^-}=-v_{P_k^-}$.  \textbf{Killed.}
\end{itemize}

\subsubsection*{Case (II.a):  $e_1$ in column $1$ at row $r_1>a$.}

\begin{itemize}[leftmargin=2em,topsep=0.3em,itemsep=0.2em]
\item In $P_k^+$: $e_1$ in column $1$ (row $r_1>a$), $e_2$ in column
$2$ (row $s_2\le a$).  Different rows and columns, axial distance
$=1+r_1-s_2\ge 2$.  \textbf{Expanded.}
\item In $P_k^-$: $e_1$ in column $1$, $e_2$ in column $1$.  Same
column, consecutive (no integer between), so consecutive rows of
column $1$.  \textbf{Killed.}
\end{itemize}

\subsubsection*{Case (II.b): $e_1$ in column $2$ at row $s_1\le a$.}

\begin{itemize}[leftmargin=2em,topsep=0.3em,itemsep=0.2em]
\item In $P_k^+$: $e_1,e_2$ both in column $2$.  Consecutive
integers, so consecutive rows of column $2$.  \textbf{Killed.}
\item In $P_k^-$: $e_1$ in column $2$ (row $s_1\le a$), $e_2$ in
column $1$ (row $r_3>a$).  Different rows and columns.  Axial
distance $|c(e_2)-c(e_1)|=|(1-r_3)-(2-s_1)|=|s_1-1-r_3|=r_3+1-s_1\ge 2$.
Non-degenerate.  \textbf{Expanded.}
\end{itemize}

In all four sub-cases (I.a)--(II.b), exactly one of $\{P_k^+,P_k^-\}$
is annihilated by $(T_{e_1}{+}1)$, and the other yields a
non-degenerate $T_{e_1}$-2-block.  This proves
Lemma~\ref{lem:kill-expand}.  $\hfill\square$

\begin{remark}
The argument depends on:
(i) the SYT-rectangular structure of $\lambda=(2^a,1^{n-2a})$
  (every cell is in column $1$ or column $2$);
(ii) the same-row exclusion
  (Corollary~\ref{cor:no-same-row}, which uses $n\ge 3a$);
(iii) the elementary fact that consecutive integers in the same
  column of a SYT occupy consecutive rows.
No special structure of the trajectory's encoding bits
$\tau\in\{L,R\}^{a-1}$ is needed.  The case bound $n\ge 3a$ enters
once, in Corollary~\ref{cor:no-same-row}.
\end{remark}

\subsection{Why no cancellation: a comment on (P1)}\label{sec:nonzero}

We used that the Hoefsmit $+q$-eigenvector in a non-degenerate
2-block has \emph{both} seminormal coefficients nonzero.  This is the
standard formula
\[
   (T_e+1)v_T=\tfrac{1}{[d]_q}\cdot[q]_q\cdot v_T+\tfrac{[d{-}1]_q[d{+}1]_q}{[d]_q^2}\cdot v_{s_eT},
\]
where $d=c(e+1)-c(e)$ (axial distance) and $[m]_q=\frac{q^m-1}{q-1}$.
For $|d|\ge 2$, $[d]_q$, $[d-1]_q[d+1]_q$, and $[d{+}1]_q$
are all nonzero in $\mathbb{Q}(q)$ (no root-of-unity collision); so
both coefficients of the resulting eigenvector are nonzero.
This gives (P1) and (P4) of Definition~\ref{def:invariant}.

\section{Consequences}\label{sec:consequences}

\subsection{Rank-zero theorem at general $a$}

\begin{theorem}[Rank-zero for $(2^a,1^{n-2a})$, $a\ge 1$,
$n\ge 3a$]\label{thm:rank-zero}
For every $a\ge 1$ and $n\ge 3a$:
$\Pi^{S_n}|_{V_{(2^a,1^{n-2a})}}=0$, conditional on the
general-$a$ Phase~A endpoint at $\lambda$
(May-14 Proposition~3.1, generalized).
\end{theorem}

\begin{proof}
Induction on $a$.  Base $a=1$: hook case, May-11
(\texttt{2026-05-11-j-formula.tex}).

Inductive step $a-1\to a$ for $a\ge 2$:  By the inductive
hypothesis at $\mu_1=(2^{a-1},1^{n-2a})$ on $n-2$ letters,
$B_{n-a}^{(\mu_1)}V_{\mu_1}\subseteq E_1^-|_{V_{\mu_1}}$.  By
May-15-gen Theorem~3.2 \eqref{eq:gen-reduction},
\[
   B_{n-a+1}^{(\lambda)}V_\lambda
   =\prod_{j=n-a}^{n-2}(T_j{+}1)\,\Phi_1\big(B_{n-a}^{(\mu_1)}V_{\mu_1}\big).
\]
$\Phi_1$ preserves $E_1^-$ (May-15-a3 Lemma~3.5 / May-15-gen
Lemma~3.3, equivariance under $T_1,\ldots,T_{n-3}$).
Each $(T_j{+}1)$ for $j\ge n-a$ has index gap from $T_1$ at least
$n-a-1\ge 2a-1\ge 3$, so commutes with $T_1$; hence preserves
$E_1^-$.  Therefore
$B_{n-a+1}^{(\lambda)}V_\lambda\subseteq E_1^-|_{V_\lambda}$.

The standard factorization (May-13 \S2 / May-15-gen Corollary~3.3)
$\Pi^{S_n}=R'_2\cdots R'_{n-a}\cdot B_{n-a+1}^{(\lambda)}$ has
leftmost factor $R'_{n-a}$ whose rightmost-in-product factor is
$(T_1{+}1)$, annihilating $E_1^-$.  So
$\Pi^{S_n}V_\lambda=0$.
\end{proof}

\begin{remark}[Status of Phase~A]
Theorem~\ref{thm:rank-zero} is conditional on Phase~A at general $a$.
Phase~A is the assertion $R'_nV_\lambda=E_{n-1}^+|_{V_\lambda}$,
which is what allows the May-15-gen Theorem~3.2 reduction
to start from $\Phi_1(V_{\mu_1})$.  This is proved rigorously at
$a=2$ (May-13 Proposition~3.1), $a=3$ (May-14 Proposition~3.1).
For $a\ge 4$ it is asserted to be a mechanical generalization
(May-15-gen Theorem~2.1) and routinely verifiable computationally
but not formalized.  Removing this conditional is a separate
project, structurally orthogonal to the per-trajectory framework.
\end{remark}

\subsection{Dimension bound: $\dim B_{n-a+1}^{(\lambda)}V_\lambda\le 1$}

Theorem~\ref{thm:main} implies $w_{a-1}(T')$ has $2$-element support
for each $T'$.  Linearly combining over $T'\in\mathrm{supp}(u^{(\mu_1)})$
(where $u^{(\mu_1)}$ generates $B_{n-a}^{(\mu_1)}V_{\mu_1}$, $1$-dim
by the inductive hypothesis) gives:
\[
   B_{n-a+1}^{(\lambda)}V_\lambda
   =\mathbb{Q}(q)\cdot\sum_{T'\in\mathrm{supp}(u^{(\mu_1)})}
     c_{T'}\,w_{a-1}(T').
\]
This is at most $1$-dimensional (a single scalar multiple of the
sum), independent of the support overlap structure.  Combined with
the standard fact $B_{n-a+1}^{(\lambda)}V_\lambda\ne 0$
(generic-$q$ argument), we conclude $\dim=1$.

\subsection{Toward the $S_a^{(n)}$ support identification}\label{sec:remaining}

Two pieces remain to identify the support as $S_a^{(n)}$ exactly
($|S_a^{(n)}|=2^a$):

\paragraph{Encoding identification.}  For each
$T'\in S_{a-1}^{(n-2)}$, the pair
$\{P_{a-1}^+(T'),P_{a-1}^-(T')\}$ equals $\{T_{\tau L},T_{\tau R}\}$
where $\tau=\varepsilon(T')$ is the $S_{a-1}$-encoding.
(May-16 Lemma~4.3.)

\paragraph{Disjointness.}  Different $T'_1,T'_2\in S_{a-1}^{(n-2)}$
give disjoint $w_{a-1}$-supports.  (May-16 Lemma~4.2.)

A clean route: track the col-$2$ entry set $\sigma(P_k^+)\cap
\{n-2a+1,\ldots,n\}$ at each step.  At step $0$,
$\sigma(P_0^+)=\sigma(T_A(T'))=\sigma(T')\cup\{n-1\}$.  Each
kill/expand step alters $\sigma$ in one position
($e_2$ or $e_3$ swaps between columns).  The recursion of $S_a$
($\{0\}\cup\mathrm{shift}(S_{a-1})\sqcup\mathrm{shift}(S_{a-1})\cup\{2a-1\}$)
matches the shape of these alterations; verifying the match
amounts to a finite induction on $a$.  This residue is not
addressed in this paper but is finite combinatorics.

\section{Verification against May-16}\label{sec:verify}

We re-derive the kill/expand pattern at $a=4$, step $1$, from
Lemma~\ref{lem:kill-expand}, and check against May-16
\S6.2 (Step~1 case analysis).

\begin{itemize}[leftmargin=2em,topsep=0.3em,itemsep=0.2em]
\item $a=4$, $n\ge 12$, $T'\in S_3^{(n-2)}$ on $n-2$ letters.
\item Step $0$: pair $\{T_A,T_B\}$, with $e_2=n-1$ at $(4,2)$,
  $e_3=n$ at $(n-4,1)$.  Case~(II) (since $e_2$ in column $2$).
\item Step $1$: apply $(T_{n-2}{+}1)$, $e_1=n-2$.
  Whether $e_1$ is in column $1$ or $2$ of $T_A$ is determined by
  $T'$: $n-2\in\sigma(T_A)$ iff $n-2\in\sigma(T')$, which holds
  iff $\tau_1=R$ in $T'$'s encoding.
\item Case~(II.a):  $e_1$ in column $1$ ($\tau_1=L$).  $P_0^+=T_A$
  expanded; $P_0^-=T_B$ killed.  Matches May-16 ``$\tau_1=L$'' row.
\item Case~(II.b):  $e_1$ in column $2$ ($\tau_1=R$).  $P_0^+=T_A$
  killed; $P_0^-=T_B$ expanded.  Matches May-16 ``$\tau_1=R$'' row.
\end{itemize}

For Step $2$ of May-16 Table~2: after Step $1$, the new pair is a
$T_{n-2}$-2-block.  In the surviving pair, $e_2^{(\text{new})}=n-2$
and $e_3^{(\text{new})}=n-1$ swap columns; the pair is in Case~(I)
or~(II) depending on the column of $n-2$.  Then Step $2$
($e_1^{(\text{new})}=n-3$) splits into sub-cases (a)/(b)
depending on the column of $n-3$ in the surviving pair, which is
determined by $\tau_2$.  All $8$ rows of Table~2 are
re-derivable as Case~(I.a)/(I.b)/(II.a)/(II.b).

(We have explicitly checked $\tau\in\{LLL,LRL\}$ against the table,
matching the kill/expand pattern in both cases.)

\section{Honest assessment}\label{sec:honest}

\paragraph{What this paper rigorously proves.}
\begin{enumerate}[leftmargin=2em,topsep=0.3em,itemsep=0.2em]
\item Theorem~\ref{thm:main} (uniform per-trajectory invariant)
for every $a\ge 2$ and $n\ge 3a$, via a four-case structural
argument.
\item Theorem~\ref{thm:rank-zero} (rank-zero theorem at
$\lambda=(2^a,1^{n-2a})$ for all $a\ge 2$, $n\ge 3a$), conditional
on general-$a$ Phase~A endpoint.
\end{enumerate}

\paragraph{What remains conditional.}
\begin{enumerate}[leftmargin=2em,topsep=0.3em,itemsep=0.2em]
\item Phase~A endpoint at general $\lambda=(2^a,1^{n-2a})$
(unchanged status from May-15-gen Remark~2.2).  Proved at
$a=2,3$; asserted ``mechanical'' for $a\ge 4$.
\item The L/R encoding identification (May-16 Lemma~4.3 generalized
to $a$): each $T'\in S_{a-1}^{(n-2)}$ produces the pair
$\{T_{\tau L},T_{\tau R}\}$ with first $a{-}1$ encoding bits matching
$\tau=\varepsilon(T')$.  This requires tracking the col-$2$ entry
set $\sigma(P_k^+)$ across the trajectory, which I have not
formalized in general.  It is the finite combinatorial residue of
the conjecture~\ref{cor:Sa-conj} (below).
\end{enumerate}

\begin{conjecture}[Support = $S_a^{(n)}$]\label{cor:Sa-conj}
For every $a\ge 2$ and $n\ge 3a$, $B_{n-a+1}^{(\lambda)}V_\lambda$
is exactly $1$-dimensional with seminormal support $S_a^{(n)}$
($|S_a^{(n)}|=2^a$), all coefficients nonzero in $\mathbb{Q}(q)$.
\end{conjecture}

\paragraph{Comparison with May-16 §6.}  May-16's path to general
$a$ identifies the same residue (encoding identification) as the
remaining substantive piece.  This paper closes the
``exactly-one-killed-per-step'' content of Claim~(1) of May-16
\S5.3 uniformly: no need for $2^{a-1}$ cases per step at each $a$.

\paragraph{Big picture.}  The per-trajectory framework now has a
\emph{uniform} foundation.  The structural lemma
(Lemma~\ref{lem:kill-expand}) explains \emph{why} the case
enumeration at each $a$ in May-13, May-15-a3, May-16 always
``works out'': it is forced by the SYT-rectangular structure of
$\lambda$ plus the column-row bound (Lemma~\ref{lem:bound}).
The encoding identification — which is the source of the
specific recursive set $S_a$ — is now isolated as a separate
combinatorial residue, structurally independent of the
``support has size $2$ at each step'' content.

\end{document}
